\subsection{A positive local repair and an indecomposable neighborhood}
The final research round supplied a positive result within this frontier.
Let $B$ be the three-color marginal map on the 90 interval weights and $H$ the five harmonic moment map. An explicit rational $90\times5$ matrix $Q$ satisfies
\[
BQ=0,\qquad HQ=I_5,\qquad
\max_v\sum_j|Q_{vj}|=1.
\]
Thus a color coupling whose weights are at least $m>0$ admits any harmonic correction $\Delta$ with $\|\Delta\|_\infty\le m$: replace its weights by $w+Q\Delta$. Color marginals stay fixed and all weights remain nonnegative.

In contrast, ordinary six-state trades that fix one anchor of each color and permute three other labels preserve all five harmonic moments. Their span has rank 47, the full kernel of the interval map. They cannot repair a midpoint discrepancy. The larger trades in $Q$ are essential for this construction.

\begin{theorem}\label{local222}
Let
\[
U_0=\frac1{8389}
\begin{pmatrix}
863&3348&-1506&984&1266&7320\\
-2664&1223&-2700&-2046&6882&-1716\\
-1986&1044&-199&7884&936&-1506\\
5496&-4866&-2592&1727&2580&366\\
-5010&-5730&1044&120&359&3348\\
-1812&-264&-7278&-258&-3708&-487
\end{pmatrix}.
\]
For every complex unitary $U$ satisfying
\[
\|U-U_0\|_{\mathrm{op}}\le\frac1{23083105288},
\]
MOC holds for arbitrary $a\in\C^6$ and
$b=(b_1,b_1,b_2,b_2,b_0,b_0)$, with arbitrary complex $b_0,b_1,b_2$.
Other placements of the three pairs use the corresponding column permutation of the center.
\end{theorem}
\begin{proof}
The exact certificate verifies $U_0U_0^T=I$ and supplies 90 strictly positive interval weights with minimum
$m=3/70375321$. Their target is
\[
p_X=|\det U[X,12]|^2,\quad
r_Z=|\det U[Z,123]|^2+|\det U[Z,124]|^2,\quad
q_Y=|\det U[Y,1234]|^2.
\]
These comprise 50 coordinates. For every actual unitary, the six intensity relations
\[
\sum_{Z\ni i}r_Z=\sum_{X\ni i}p_X+\sum_{Y\ni i}q_Y
\]
and the mass relation $\sum_Zr_Z=2\sum_Xp_X$ hold. They follow from the projection intensities $P_{ii}$, $2P_{ii}+W_{ii}$ and $(P+W)_{ii}$, where $P$ uses columns $12$ and $W$ columns $34$. The seven relations are independent, and the interval map has rank 43, so every actual target belongs to its column space.

The certificate selects 43 rows and 43 interval columns forming an invertible matrix with inverse infinity operator norm 82. Any change of the full target in that column space is therefore realized by changing those 43 weights, with magnitude at most 82 times the target's infinity-norm change.

For two unitaries, telescoping the $k$ determinant columns and applying Hadamard's inequality gives
\[
\left||\det U[I,J]|^2-|\det V[I,J]|^2\right|
\le2k\|U-V\|_{\mathrm{op}}.
\]
The target coordinates change by at most $12\|U-U_0\|_{\mathrm{op}}$. The radius is exactly $m/(12\cdot82)$, so repaired weights stay nonnegative. All 50 equations hold because both targets lie in the same rank-43 space.

The interval law represents the canonical product input with column classes $12$, $34$, $56$ on three projective lines. A common $GL_2$ transformation transports it to any three distinct projective input lines; scalar pairing proves the spectral conclusion. Coincident values follow by continuity.
\end{proof}

The first two-column projector of $U_0$ has no zero off-diagonal entry, excluding a nontrivial reducing coordinate partition. Thus this is a certified neighborhood beyond the independent-coordinate-block class. The radius is conservative. The universal $(2,2,2)$ statement for arbitrary $U(6)$ remains unproved.
